\documentclass{article}
\usepackage{graphicx}
\usepackage{amsfonts}
\usepackage{amsmath}
\usepackage{amsthm}
\usepackage{tikz}

\theoremstyle{definition}
\newtheorem{theorem}{Theorem}[section]
\newtheorem{corollary}[theorem]{Corollary}
\newtheorem{proposition}[theorem]{Proposition}
\newtheorem{lemma}[theorem]{Lemma}
\newtheorem{definition}[theorem]{Definition}

\theoremstyle{remark}
\newtheorem{remark}[theorem]{Remark}
\newtheorem{example}[theorem]{Example}
\newtheorem{claim}[theorem]{Claim}

\title{Integral and the Fundamental Theorem of Calculus}
\author{Sarah Liu}
\date{January 2026}

\begin{document}

\maketitle

\section{Basics}
\begin{definition}
    Given a function $f(N)$,
    \begin{enumerate}
        \item $O(f(N))$ denotes the set of all $g(N)$ such that $|g(N)/f(N)|$ is bounded above as $N\rightarrow\infty$
        \item $\Omega(f(N))$ denotes the set of all $g(N)$ such that $|g(N)/f(N)|$ is bounded from below by a (strictly) positive number as $N\rightarrow\infty$
        \item $\Theta(f(N))$ denotes the set of all $g(N)$ such that $|g(N)/f(N)|$ is bounded from both above and below as $N\rightarrow\infty$
    \end{enumerate}
\end{definition}

\end{document}